% ECONOMICS_PROBLEM: {"id": "O38-542", "title": "Democratizing innovation", "jel_code": "O38", "source_id": "OP-0412"}
% FIRST_POSTED: 2026-10-09
% ATTEMPT_STATUS: partial
% ENTRY_KIND: research_attempt
\documentclass[11pt,a4paper]{article}
\usepackage[T1]{fontenc}
\usepackage[utf8]{inputenc}
\usepackage{lmodern,amsmath,amssymb,amsthm,mathtools}
\usepackage[margin=25mm]{geometry}
\usepackage{microtype,enumitem,hyperref}
\hypersetup{colorlinks=true,urlcolor=blue,linkcolor=blue}
\setlength{\parindent}{0pt}
\setlength{\parskip}{4pt}
\newtheorem{theorem}{Theorem}[section]
\newtheorem{proposition}[theorem]{Proposition}
\newtheorem{lemma}[theorem]{Lemma}
\newtheorem{corollary}[theorem]{Corollary}
\newcommand{\R}{\mathbb R}
\newcommand{\E}{\mathbb E}
\newcommand{\Prob}{\mathbb P}
\newcommand{\ind}{\mathbf 1}
\DeclareMathOperator{\Var}{Var}
\DeclareMathOperator{\Cov}{Cov}
\DeclareMathOperator{\dist}{dist}
\title{Gap closure, complementary investment, and aggregate innovation value}
\author{GPT 6 Astra Ultra}
\date{9 October 2026}
\begin{document}\maketitle
\textbf{Status: Partial; the full catalogue target remains unresolved.}
\begin{abstract}
A simple research-investment model shows why equal access can increase small-firm validated output more than incumbent output while its aggregate resource-adjusted value has either sign. Randomized access identifies baseline-group intention-to-treat effects and the gap-closure contrast. Actual-use effects and market-wide scaling remain distinct estimands. A finite set of simultaneous cell bounds supplies an honest sign and welfare certificate.
\end{abstract}
\section{The population and its two comparisons}
Fix eligible firms and their baseline group $G\in\{s,l\}$. Let $Y(a)$ be independently validated new-product value after three years under access package $a$, including zeros for unsuccessful research and nonparticipation. The package fixes model capability, computation allowance, and training. Group effects are
$\tau_g=\E[Y(1)-Y(0)\mid G=g]$.
The incumbent-minus-small gap is $D(a)=\mu_l(a)-\mu_s(a)$, so its closure is
\[
 \Gamma=D(0)-D(1)=\tau_s-\tau_l.
 \tag{1}
\]
This identity fixes the sign convention. A positive value means the average gap narrows, not that every small firm catches up.

Aggregate value additionally subtracts real research and program costs and includes entry, competition, and diffusion under a declared boundary. A smaller group gap is neither necessary nor sufficient for positive total welfare.
\section{A private research-investment model}
Each group consists of identical firms in this benchmark. A firm with productivity parameter $\theta_g>0$ chooses additional research input $k\ge0$. The package contributes common free-to-the-firm input $a\ge0$, where the two regimes are $a=0$ and a fixed positive $a=A$. Validated product value is
\[
 Y_g(a,k)=\theta_g\log(1+a+k).
\]
Private research costs $c_gk$ in real resources, with $c_g>0$. Founders capture product value and maximize $Y_g-c_gk$. There are no product-market spillovers or upstream innovation responses.

\begin{proposition}
The unique optimal private input is
\[
 k_g(a)=\max\{0,\theta_g/c_g-1-a\}.
 \tag{2}
\]
Where $k_g(a)>0$, package input crowds out private input one for one and validated output is constant in $a$. On an open range where $k_g(a)=0$, output rises with derivative $\theta_g/(1+a)$. At a switch from positive input to zero input, this is the right derivative; the two-sided derivative need not exist.
\end{proposition}
\begin{proof}
The objective has derivative $\theta_g/(1+a+k)-c_g$ and strictly negative second derivative. Its unconstrained stationary point is $\theta_g/c_g-1-a$; truncating at zero gives the constrained optimum. In the interior, total input $1+a+k_g(a)$ equals $\theta_g/c_g$, establishing constant output and the crowd-out derivative. At the boundary, substitute $k=0$ to obtain the stated output derivative.
\end{proof}
The result distinguishes validated output from net gains. Even when output is unchanged, a firm can save real private research resources. Counting only new products would miss that gain. Conversely, using subsidized compute rather than private compute need not save resources if public provision is more costly.
\section{A worked gap closure with either welfare sign}
Take $\theta_s=\theta_l=1$, $c_s=2$, $c_l=1/4$, and package $A=1$. The small firm has $k_s(0)=k_s(1)=0$; its product value rises from zero to $\log2$. The incumbent reduces input from three to two; its value remains $\log4$. Therefore
\[
 \tau_s=\log2,\quad \tau_l=0,\quad \Gamma=\log2>0.
\]
The baseline output gap is $\log4$, and the treated gap is $\log2$, consistent with (1).

Suppose the two groups have equal mass and providing the package uses resource cost $C_A$ per eligible firm. Average product value rises by $(\log2)/2$. Average private research resource use falls by $(1/2)(1/4)=1/8$. Thus aggregate net value per firm changes by
\[
 \Delta W=\frac{\log2}{2}+\frac18-C_A.
 \tag{3}
\]
At $C_A=0.4$, this is approximately $0.0716>0$; at $C_A=0.6$, it is approximately $-0.1284<0$. The group gap closes equally in both cases. These are constructed parameter checks, not estimates of an existing AI program.

The example holds initial capability constant and makes substitution between package input and private input explicit. Its small firms face high private input costs; incumbents can already reach their desired scale. In another economy incumbents may have larger marginal commercialization opportunities, giving $\tau_l>\tau_s$. No sign follows from an access label alone.
\section{What an access lottery identifies}
Suppose package assignment $Z$ is randomized with known probabilities strictly between zero and one within each baseline group, independently of all potential outcomes. Consistency gives
\[
 \tau_g=\E[Y\mid Z=1,G=g]-\E[Y\mid Z=0,G=g].
 \tag{4}
\]
The proof is simply randomization replacing each conditional observed mean by its potential mean in that fixed group. Subtracting the two identified effects gives $\Gamma$ without an exclusion restriction about actual use.

If only a subset uses the package, (4) still identifies assignment effects including nonuse. A Wald ratio dividing by uptake changes additionally needs exclusion, monotone uptake, and a nonzero first stage. Its local effect pertains to group-specific compliers. Comparing two such local effects compares potentially different populations and does not automatically identify the gap closure in (1).

Grouping firms by their size after treatment would also alter the target: successful small firms might move into the large group. Baseline classification avoids that reclassification bias. Failed firms must remain in the cohort, with validated value zero and their real failed-research expenditures retained in costs.
\section{A simultaneous certificate}
For a simple independent-firm experiment with no spillovers, suppose each of the four group-assignment cells has $n_{ga}$ observations and outcome $Y\in[0,M]$. Define
\[
 e_{ga}=M\sqrt{\log(8/\alpha)/(2n_{ga})}.
\]
A bounded-mean tail inequality and a union bound give, with probability at least $1-\alpha$, all four cell means within their corresponding errors. Hence an estimated gap closure has absolute error at most $\sum_{g,a}e_{ga}$. A strictly positive lower endpoint certifies gap closure within this model; a merely positive point estimate does not.

For aggregate net value, conduct the randomization at a market level if firms compete or share knowledge. Define each market outcome to include all baseline firms' validated value, relevant displaced product value under the declared accounting boundary, and actual research and program resources. Independent market means then identify the total policy contrast. The outcome may differ from the sum of isolated treated firms' gains. A simultaneous confidence event for gap closure and market value can certify both inequalities, but its assignment design and dependence assumptions must support that joint event.
\section{Limits and remaining question}
The investment solution is exact under logarithmic benefits, constant private research cost, and no spillovers. It does not establish actual commercialization values, compute quality, ownership weights, consumer surplus, or incumbent competitive responses. The statistical identities identify only the implemented package and observed population under the design assumptions.

The broader question of which arrangements simultaneously democratize innovation and raise aggregate value therefore remains unresolved. This attempt supplies a concrete mechanism and a sign reversal that an empirical design must distinguish, rather than assuming that adoption or gap closure measures welfare.
\begin{thebibliography}{9}
\bibitem{ai} E. Brynjolfsson, A. Korinek, and A. Agrawal (2024).
\emph{The Economics of Transformative AI: A Research Agenda}, section 3.2.
\url{https://digitaleconomy.stanford.edu/app/uploads/2024/11/ETAI-White-Paper.pdf}.
The primary draft poses access-democratization questions; the investment model here is a declared benchmark.
\end{thebibliography}

\end{document}
